\documentclass[12pt,a4paper]{report}

\usepackage[T1]{fontenc}
\usepackage[utf8]{inputenc}
\usepackage[english]{babel}
\usepackage{lmodern}
\usepackage{microtype}
\usepackage[a4paper,margin=1in]{geometry}
\usepackage{amsmath,amssymb,amsthm,mathtools}
\usepackage{enumitem}
\usepackage{hyperref}
\usepackage{csquotes}
\usepackage[backend=biber,natbib=true]{biblatex}
\addbibresource{references.bib}
\usepackage{setspace}
\usepackage{parskip}

\hypersetup{
  colorlinks=true,
  linkcolor=blue,
  citecolor=blue,
  urlcolor=blue,
  pdftitle={Toward a Probabilistic Semantics for the Lambek Calculus},
  pdfauthor={}
}

\onehalfspacing

\newtheorem{theorem}{Theorem}[chapter]
\newtheorem{proposition}[theorem]{Proposition}
\newtheorem{lemma}[theorem]{Lemma}
\newtheorem{corollary}[theorem]{Corollary}
\theoremstyle{definition}
\newtheorem{definition}[theorem]{Definition}
\newtheorem{remark}[theorem]{Remark}
\newtheorem{example}[theorem]{Example}

\newcommand{\LC}{\mathsf{L}}
\newcommand{\At}{\mathsf{At}}
\newcommand{\Frag}{\mathsf{Frag}}
\newcommand{\Lex}{\mathsf{Lex}}
\newcommand{\Sem}{\models_{\mathcal C}}
\newcommand{\concat}{\mathbin{\cdot}}
\newcommand{\Pow}{\mathcal{P}}

\title{Toward a Probabilistic Semantics for the Lambek Calculus\\\large Formalisation, Soundness, and Directions for Completeness}
\author{[Author Name]}
\date{[Month Year]}

\begin{document}

\maketitle
\pagenumbering{roman}

\begin{abstract}
This thesis studies a corpus-relative probabilistic semantics for the associative Lambek
calculus. The Lambek calculus was introduced by Lambek as a logic of syntactic types and has
since become one of the foundational systems of categorial and type-logical grammar.
Its appeal lies in the tight match between syntax and composition: grammatical derivations
track how word meanings combine, which makes the calculus attractive both in formal
linguistics and in mathematically oriented approaches to natural language processing.

The central motivation of this project is ambiguity. In a purely logical parser, a sentence may
admit several derivations, all equally well-formed at the level of derivability. In practical
language processing, however, one usually wants not merely to decide grammaticality but also
to rank candidate analyses. Probabilistic grammars such as PCFGs and probabilistic CCGs do
exactly this, and their empirical success suggests that a probabilistic version of Lambek-style
parsing is worth developing. The aim here is therefore to define a semantics that assigns to a
string and a Lambek type a value in $[0,1]$, interpreted as the probability or confidence that
the string realises that type relative to a finite corpus.

The thesis contributes a cleaned-up formalisation of this idea. The semantics is restricted to a
finite set of observed corpus fragments, which keeps the recursive clauses well-defined and makes
explicit the sense in which the model is corpus-relative. Particular attention is paid to technical
subtleties that obstructed the earlier draft: words outside the corpus vocabulary are excluded
from the formal domain instead of being given ad hoc values, atomic types not attested for a
word receive probability $0$, and the residual clauses are defined so as to avoid division by
zero. On the basis of these definitions, a soundness theorem is proved for the algebraic axioms
of the associative Lambek calculus.

The thesis does \emph{not} claim completeness for the probabilistic semantics introduced here.
On the contrary, it argues that completeness is the natural next problem. Classical completeness
results for the Lambek calculus are already known for several non-probabilistic semantics, but
it remains open whether an analogous theorem can be obtained for the present corpus-based
probabilistic setting. The final chapter therefore outlines future work on completeness,
implementation, smoothing for unseen data, and empirical evaluation against existing
probabilistic parsing frameworks.
\end{abstract}

\tableofcontents
\clearpage
\pagenumbering{arabic}

\chapter{Introduction}

\section{Why the Lambek calculus matters}
The Lambek calculus was introduced in \citet{Lambek1958} as a formal calculus of syntactic
types. Its original purpose was already recognisably linguistic: to give a principled method for
distinguishing well-formed from ill-formed strings and to model the way in which syntactic
categories combine in natural language. Lambek's follow-up work in \citet{Lambek1961}
further developed the type-theoretic viewpoint and helped establish categorial grammar as a
serious mathematical approach to syntax.

The importance of the Lambek calculus lies in the fact that it sits at an unusually fruitful
intersection of logic, algebra, and linguistics. On the logical side, it is one of the prototypical
substructural logics: order matters, and structural principles such as exchange, weakening, and
contraction are not built in by default. On the linguistic side, this resource sensitivity is not an
artificial restriction but an advantage, because word order and local composition are central to
syntactic analysis. For this reason the Lambek calculus became a foundational system in the
study of categorial and type-logical grammars \citep{vanBenthem1988,Moortgat1995Multimodal,morrill1996generalising}.

A second reason for its importance is that the calculus offers an explicit bridge between syntax
and semantics. Derivations in categorial calculi naturally correspond to compositional analyses of
meaning, and this connection has been one of the main reasons for the enduring influence of
Lambek-style systems in formal semantics and mathematical linguistics \citep{vanBenthem1988,Steedman2000}.

\section{Completeness: classical results and the open problem of this thesis}
There is an important distinction to make at the outset. The completeness of the \emph{classical}
Lambek calculus is not an open problem in general. Several major completeness results are
already known. \citet{Dosen1985} proved a completeness theorem for the Lambek calculus of
syntactic categories; \citet{AndrekaMikulas1994} established strong completeness results for
relational semantics; and \citet{Pentus1995} proved completeness with respect to language
models. In addition, Pentus' work on the context-free status of Lambek grammars clarified the
place of the calculus in formal-language theory \citep{Pentus1993,Pentus1997}.

The open problem relevant to the present thesis is therefore narrower and more specific:
\emph{is the probabilistic semantics introduced here complete for the Lambek calculus?} This
question is open because the semantics developed in this thesis is corpus-relative, numerical,
and intentionally partial in the sense that it is tied to observed linguistic data. The standard
canonical-model and algebraic techniques from classical completeness proofs do not transfer
immediately to such a setting.

For that reason the present thesis focuses on soundness first. Establishing soundness is the right
first milestone: before asking whether the semantics captures \emph{all} valid inferences, one
must first show that every derivable Lambek inequality is semantically valid. Completeness is the
next goal, not a claim of the present work.

\section{Lambek calculus in linguistics}
In linguistics, the Lambek calculus is best understood as part of the broader tradition of
categorial grammar. Categories behave like types, and grammatical composition is governed by
logical operations rather than by phrase-structure rules. This makes the calculus attractive for
analysing phenomena where lexical information and local combinatorics are central.

Type-logical grammar developed this perspective in depth and connected it to proof theory,
formal semantics, and computational linguistics \citep{Moortgat1997,Morrill1994}. Closely related
systems, especially Combinatory Categorial Grammar (CCG), have been particularly successful in
wide-coverage parsing and semantic interpretation \citep{Steedman2000,ClarkCurran2007}.
Although CCG is not identical to the Lambek calculus, the family resemblance is close enough
that probabilistic work on CCG provides a useful source of intuition for the present project.

\section{Why probabilistic semantics?}
Pure derivability is not enough for realistic parsing. A sentence may have several syntactically
well-formed derivations, and those derivations may correspond to different readings. This is not a
pathology but a normal feature of natural language. A parser that returns every derivation with
equal status is therefore mathematically respectable but practically incomplete.

Probabilistic grammatical models address this by associating weights or probabilities with lexical
choices and structural combinations. In probabilistic context-free grammars, probabilities rank
parse trees generated by grammar rules; in probabilistic CCG, similar ideas are applied to a more
lexicalised grammar formalism \citep{ManningSchutze1999,ClarkCurran2007}. There has also been
explicit prior work on stochastic Lambek grammars \citep{BonfanteDeGroote2004}. The basic
intuition is simple: if several analyses are licensed, prefer the one better supported by the data.

This thesis adopts that intuition in a deliberately conservative way. Rather than starting from a
full statistical parser, it starts from a semantic valuation $v(x,A) \in [0,1]$ for strings $x$ and
Lambek types $A$, interpreted relative to a corpus. The main question is then whether the
algebraic behaviour of the Lambek connectives can be captured in a way that supports a
sound proof theory.

\section{Aims and scope of the thesis}
The thesis has four main aims.
\begin{enumerate}[label=(\roman*)]
    \item To motivate a probabilistic semantics for the associative Lambek calculus from the
    perspective of parsing and ambiguity resolution.
    \item To formalise that semantics carefully enough that the main recursive clauses are
    mathematically well-defined.
    \item To prove soundness of the Lambek axioms with respect to the proposed semantics.
    \item To identify the obstacles that make completeness the next major step rather than a result
    already achieved.
\end{enumerate}

Several things are \emph{not} attempted here. The thesis does not provide a learning algorithm for
extracting lexical probabilities from raw corpora, does not present a full parser implementation,
and does not prove completeness. These tasks are deferred to future work.

\section{Suggested background reading}
A useful reading path for the topics covered in this thesis is the following.
\begin{itemize}
    \item \citet{Lambek1958} and \citet{Lambek1961} for the original calculus of syntactic types.
    \item \citet{vanBenthem1988}, \citet{Moortgat1997}, and \citet{Morrill1994} for type-logical and
    categorial grammar background.
    \item \citet{AndrekaMikulas1994}, \citet{Dosen1985}, and \citet{Pentus1995} for completeness
    results and model theory.
    \item \citet{GalatosJipsenKowalskiOno2007} for the broader algebraic perspective on substructural logics.
    \item \citet{ManningSchutze1999}, \citet{Steedman2000}, and \citet{ClarkCurran2007} for
    probabilistic parsing and CCG.
    \item \citet{BonfanteDeGroote2004} for a directly related stochastic approach to Lambek
    grammars.
\end{itemize}

\chapter{Literature Review}

\section{Classical Lambek calculus and completeness}
The classical literature provides an important baseline for this thesis. The original system of
\citet{Lambek1958} and \citet{Lambek1961} was quickly recognised as more than a linguistic
notation: it is a mathematically robust substructural logic with a rich proof theory and model
theory. Surveys and handbook treatments such as \citet{vanBenthem1988} and
\citet{Moortgat1997} show how the calculus sits within the wider landscape of categorial and
type-logical grammar.

For the specific question of completeness, the paper by \citet{Dosen1985} is especially relevant
because it directly establishes a completeness theorem for the Lambek calculus of syntactic
categories. Equally important is the relational-semantics perspective of
\citet{AndrekaMikulas1994}, who showed that relational models yield strong completeness results
for suitable versions of the calculus. From a language-theoretic viewpoint, \citet{Pentus1995}
proved completeness with respect to language models, while \citet{Pentus1993} and
\citet{Pentus1997} clarified the relation between Lambek grammars and context-free grammars.

These results matter for the present thesis in two ways. First, they demonstrate that
completeness questions for Lambek systems are deep but tractable. Second, they show that
there are already several mathematically natural semantic classes for which completeness holds.
This makes it reasonable to ask whether a similar story can be told for a probabilistic,
corpus-based semantics.

\section{Substructural logic and algebraic background}
The Lambek calculus belongs to the broader family of substructural logics, where structural
rules are controlled rather than assumed. This family includes relevance logics, linear logics, and
other resource-sensitive systems. Standard algebraic treatments organise these systems around
residuated structures, in which multiplication and division are linked by residuation principles
\citep{GalatosJipsenKowalskiOno2007,DosenSchroederHeister1993}.

This background is helpful because the semantics proposed in this thesis is explicitly designed to
mirror residuation: the multiplication clause aggregates compatible decompositions of a string,
while the two division clauses are defined by corpus-relative ratios. The result is not a standard
residuated algebra, but it is clearly inspired by the same adjoint pattern.

\section{Lambek calculus in linguistics and computational semantics}
Within linguistics, Lambek-style systems are valued because they combine lexical sensitivity,
word order, and compositionality in a transparent way. Textbook and monograph treatments
such as \citet{Morrill1994} and \citet{Moortgat1997} develop the linguistic applications in detail.
Steedman's work on CCG shows how categorial methods support wide-coverage syntactic and
semantic analysis \citep{Steedman2000}.

A further line of work connects categorial grammar to compositional models of meaning. While
this thesis is not primarily about vector semantics, the compositional viewpoint remains relevant:
if syntactic derivations guide semantic composition, then a probabilistic ranking over derivations
can also be interpreted as a ranking over candidate semantic analyses. For broader motivation on
composition-driven semantics, see \citet{CoeckeSadrzadehClark2010}.

\section{Probabilistic parsing: PCFGs, probabilistic CCGs, and stochastic Lambek grammars}
The most direct computational motivation for this thesis comes from probabilistic parsing.
Probabilistic context-free grammars are the standard starting point for ranking parses in the
presence of ambiguity \citep{ManningSchutze1999}. They are attractive because they combine a
clear formal semantics with trainable numerical parameters.

Probabilistic CCG takes similar ideas into a more lexicalised grammar formalism.
\citet{ClarkCurran2007} developed a wide-coverage statistical CCG parser using log-linear models,
showing that categorial methods can be scaled successfully when combined with probabilistic
ranking. For the Lambek setting more directly, \citet{BonfanteDeGroote2004} proposed a
stochastic point of view on Lambek categorial grammars, motivated precisely by the fact that
parsing in the Lambek calculus is inherently non-deterministic.

These works suggest that a probabilistic Lambek semantics is not an isolated curiosity. It sits at
a natural intersection of logical grammar and statistical parsing.

\section{Recent model-theoretic developments}
Recent work by Stepan Kuznetsov and collaborators shows that semantics and completeness for
Lambek-style systems remain active research topics. Of particular relevance is
\citet{Kuznetsov2023}, which revisits relational semantics for extensions of the Lambek calculus
and carefully distinguishes weak from strong completeness. Likewise,
\citet{KanovichKuznetsovScedrov2022} studies language models for extensions of the Lambek
calculus and highlights how adding connectives or constants can significantly affect completeness
behaviour.

For this thesis, the lesson is methodological: completeness is subtle even in well-studied,
non-probabilistic extensions. That makes it plausible that the completeness problem for a
corpus-relative probabilistic semantics will require new ideas rather than a straightforward
adaptation of existing proofs.

\section{Position of the present thesis}
The literature therefore leaves a clear space for the present project. On one side we have a
well-developed proof theory and completeness theory for classical Lambek calculi. On the other
side we have successful probabilistic grammar formalisms, together with some directly related
stochastic Lambek work. What is still missing is a careful, thesis-length account of a
corpus-based probabilistic semantics that is explicit about its formal choices, honest about its
limitations, and rigorous at least on soundness. That is the niche this thesis aims to fill.

\chapter{Formalisation}

\section{Parsing goal and modelling stance}
The intended task is the following: given a finite corpus $\mathcal C$ of sentences and a string of
words drawn from the corpus vocabulary, we want to assign Lambek types to that string and,
when several parses are available, rank them probabilistically. In other words, the model is not
only a recogniser of grammaticality but a \emph{ranker} of candidate analyses.

The formalisation in this thesis deliberately simplifies several practical issues.
\begin{itemize}
    \item It assumes that lexical type probabilities are already available, either from manual
    annotation or from a prior learning stage.
    \item It does not handle words outside the corpus vocabulary.
    \item It treats probabilities as corpus-relative confidence values rather than as a fully learned
    generative model.
    \item It abstracts away from implementation issues such as chart parsing, estimation, and
    smoothing.
\end{itemize}
These choices are not defects so much as scope restrictions: they isolate the semantic core that
is needed for the soundness proof.

\section{Syntax of the associative Lambek calculus}
Let $\At$ be a non-empty set of atomic types. The set of Lambek formulae is generated by the
grammar
\[
A ::= p \mid A \backslash A \mid A / A \mid A \concat A
\qquad (p \in \At).
\]
We work with the associative Lambek calculus, written $\LC$, in an algebraic presentation. The
primitive relation $\vdash_{\LC}$ is a preorder on formulae generated by the following validities:
\begin{align*}
& A \vdash_{\LC} A, \\
& (A \concat B) \concat C \vdash_{\LC} A \concat (B \concat C), \\
& A \concat (B \concat C) \vdash_{\LC} (A \concat B) \concat C, \\
& A \concat B \vdash_{\LC} C \iff A \vdash_{\LC} C/B, \\
& A \concat B \vdash_{\LC} C \iff B \vdash_{\LC} A\backslash C,
\end{align*}
and transitivity.

This is equivalent in spirit to the usual residuated presentation of the associative Lambek
calculus. It is particularly convenient here because the soundness proof proceeds by checking the
validity of identity, associativity, residuation, and transitivity under the proposed semantics.

\section{Corpus fragments and vocabulary}
Let $\mathcal C$ be a finite corpus of sentences over a finite vocabulary $\Sigma$. We denote $\Sigma^+$ to be the set of finite non-empty strings from vocabulary $\Sigma$. In order to keep the
semantics finite and directly tied to parsing data, let $\Frag(\mathcal C)$ be the set of all non-empty
contiguous strings that occur in at least one sentence in $\mathcal C$. Thus:
\[
\Frag(\mathcal C) \subseteq \Sigma^+.
\]
Because $\mathcal C$ is finite, the set $\Frag(\mathcal C)$ is finite as well.

The choice of $\Frag(\mathcal C)$ rather than all of $\Sigma^+$ is deliberate. It reflects the fact that
this semantics is trained and evaluated relative to observed data. It also ensures that the minima
and maxima appearing in the recursive clauses below are taken over finite sets. Extending the
semantics to arbitrary strings over $\Sigma$ would require a separate treatment of unseen data,
which is left to future work.

\section{Lexical probabilities}
We assume a lexical probability assignment
\[
\Lex : \Sigma \times \At \to [0,1]
\]
such that for every word $w \in \Sigma$,
\[
\sum_{p \in \At} \Lex(w,p) = 1,
\]
where all but finitely many summands are zero. Intuitively, $\Lex(w,p)$ is the probability that
the word $w$ is assigned the atomic type $p$.

\begin{remark}
If a word $w \in \Sigma$ occurs in the corpus but is never assigned an atomic type $p$, then
$\Lex(w,p)=0$. This matches the intuition from the draft: lack of attestation for a type within
the observed vocabulary should count against that type.
\end{remark}

\begin{remark}
Words outside the corpus vocabulary are not assigned probabilities here at all. Rather than
assigning them an arbitrary value such as $0$ or leaving them as exceptional elements with
undefined behaviour, they are simply excluded from the domain of the semantics. This keeps the
formal system clean and makes explicit that out-of-vocabulary handling is a future extension.
\end{remark}

\section{Probabilistic valuation}
We now define a valuation
\[
v : \Frag(\mathcal C) \times \mathrm{Fm}(\LC) \to [0,1].
\]
The definition proceeds by recursion on formula complexity.

\begin{definition}[Atomic clauses]
For $x \in \Frag(\mathcal C)$ and $p \in \At$, define
\[
v(x,p) =
\begin{cases}
\Lex(x,p) & \text{if $x$ is a single word in $\Sigma$,}\\
0 & \text{if $x$ has length greater than $1$.}
\end{cases}
\]
\end{definition}

The second clause is a modelling choice: atomic categories are taken to be lexical categories of
single words. One could allow atomic categories for longer fragments as well, but that would
require extra data and is unnecessary for the present soundness result.

\begin{definition}[Product]
For $x \in \Frag(\mathcal C)$,
\[
v(x,A\concat B)
= \max\{v(y,A)\,v(z,B) \mid y,z \in \Frag(\mathcal C),\ yz=x\},
\]
with the convention that the maximum of the empty set is $0$.
\end{definition}

This says that the probability of $x$ having composite type $A\concat B$ is the best score among
all observed binary decompositions of $x$ into a left part of type $A$ and a right part of type
$B$.

\begin{definition}[Left residual]
For $x \in \Frag(\mathcal C)$,
\[
v(x,A\backslash B)
= \min\Bigl(\{1\} \cup
\Bigl\{\min\bigl(1,\frac{v(w,B)}{v(y,A)}\bigr)
\ \Bigm|\ y \in \Frag(\mathcal C),\ yx=w \in \Frag(\mathcal C),\ v(y,A)>0\Bigr\}\Bigr).
\]
\end{definition}

\begin{definition}[Right residual]
For $x \in \Frag(\mathcal C)$,
\[
v(x,B/A)
= \min\Bigl(\{1\} \cup
\Bigl\{\min\bigl(1,\frac{v(w,B)}{v(y,A)}\bigr)
\ \Bigm|\ y \in \Frag(\mathcal C),\ xy=w \in \Frag(\mathcal C),\ v(y,A)>0\Bigr\}\Bigr).
\]
\end{definition}

\section{Why these clauses are well-defined}
The original draft raised two technical problems: division by zero in the residual definitions and
the treatment of words not present in the corpus. The present formalisation resolves both.

First, division by zero is avoided by restricting the residual minima to witnesses with positive
denominator, i.e. to contexts $y$ such that $v(y,A)>0$. If no such witness exists, the set in the
minimum collapses to $\{1\}$ and the value becomes $1$. This is a natural default: in the absence
of any positively supported counterevidence, the residual connective is not forced below $1$.
The inner truncation $\min(1,\_)$ ensures that the value stays in the interval $[0,1]$.

Second, out-of-vocabulary words are excluded from the domain entirely. The semantics is not a
universal language model over all possible strings; it is a finite, corpus-relative semantics for
observed fragments. This is consistent with the stated goal of the thesis, namely to parse strings
of words drawn from the corpus vocabulary.

\begin{remark}[Subtleties not captured by the present formalisation]
The following issues are intentionally left outside the formal model.
\begin{enumerate}[label=(\alph*)]
    \item \textbf{Learning from raw corpora.} The formalisation assumes lexical probabilities are
    already given, whereas a practical system would need to infer them.
    \item \textbf{Out-of-vocabulary generalisation.} Smoothing, backoff, or subword modelling would be
    needed for real applications.
    \item \textbf{Global normalisation.} The values $v(x,A)$ behave as probabilities locally, but the model is
    not presented here as a full generative probability distribution over derivation trees.
    \item \textbf{Unseen but grammatical strings.} By working over $\Frag(\mathcal C)$, the semantics is tied to
    attested fragments rather than all strings over $\Sigma$.
\end{enumerate}
These restrictions should be read as part of the thesis scope, not as permanent limitations on
future versions of the model.
\end{remark}

\section{Semantic consequence}
\begin{definition}
For formulae $A$ and $B$, define
\[
A \Sem B
\quad\text{iff}\quad
v(x,A) \leq v(x,B) \text{ for all } x \in \Frag(\mathcal C).
\]
\end{definition}

Thus $A \Sem B$ means that every observed fragment is at least as probable as a realisation of
$B$ as it is as a realisation of $A$. This is the semantic preorder with respect to which soundness
will be proved.

\chapter{Soundness}

\section{Strategy of proof}
We now prove that the proposed semantics is sound for the associative Lambek calculus. Since
$\LC$ is generated by identity, associativity, residuation, and transitivity, it is enough to verify
that each of these principles is valid under $\Sem$.

\begin{theorem}[Identity]
For every formula $A$, we have $A \Sem A$.
\end{theorem}

\begin{proof}
For every $x \in \Frag(\mathcal C)$, we trivially have $v(x,A) \leq v(x,A)$. Hence $A \Sem A$.
\end{proof}

\section{Associativity of product}
\begin{theorem}[Associativity]
For all formulae $A,B,C$,
\[
(A\concat B)\concat C \Sem A\concat(B\concat C)
\qquad\text{and}\qquad
A\concat(B\concat C) \Sem (A\concat B)\concat C.
\]
\end{theorem}

\begin{proof}
Let $x \in \Frag(\mathcal C)$. By definition,
\begin{align*}
v(x,(A\concat B)\concat C)
&= \max\{v(u,A\concat B)\,v(w,C) \mid uw=x\} \\
&= \max\{v(y,A)\,v(z,B)\,v(w,C) \mid yz=u,\ uw=x\}.
\end{align*}
Since string concatenation is associative, choosing $u,w$ with $uw=x$ and then $y,z$ with
$yz=u$ is equivalent to choosing a triple $y,z,w$ with $yzw=x$. Therefore
\[
v(x,(A\concat B)\concat C)
= \max\{v(y,A)\,v(z,B)\,v(w,C) \mid yzw=x\}.
\]
Running the same calculation from the other bracketing gives
\[
v(x,A\concat(B\concat C))
= \max\{v(y,A)\,v(z,B)\,v(w,C) \mid yzw=x\}.
\]
Hence the two values are equal for every $x$, and the first semantic inequality follows.
The converse direction is identical.
\end{proof}

\section{Residuation}
The key step is to show that product and division satisfy the expected adjoint behaviour.

\begin{proposition}[Right residuation]
For all formulae $A,B,C$,
\[
A\concat B \Sem C
\quad\Longleftrightarrow\quad
A \Sem C/B.
\]
\end{proposition}

\begin{proof}
Suppose first that $A\concat B \Sem C$. Let $x \in \Frag(\mathcal C)$. We must show that
$v(x,A) \leq v(x,C/B)$.

If the defining set for $v(x,C/B)$ contributes no witness beyond the default element $1$, then
$v(x,C/B)=1$ and the claim is immediate because $v(x,A) \in [0,1]$.

Otherwise take any $y \in \Frag(\mathcal C)$ such that $xy=w \in \Frag(\mathcal C)$ and
$v(y,B)>0$. Since $xy=w$, the decomposition $w=xy$ is one of the candidates in the product
clause, hence
\[
v(w,A\concat B) \geq v(x,A)\,v(y,B).
\]
By the assumption $A\concat B \Sem C$, we have $v(w,A\concat B) \leq v(w,C)$. Therefore
\[
v(x,A)\,v(y,B) \leq v(w,C).
\]
Because $v(y,B)>0$, division yields
\[
v(x,A) \leq \frac{v(w,C)}{v(y,B)}.
\]
As also $v(x,A)\leq 1$, we obtain
\[
v(x,A) \leq \min\Bigl(1,\frac{v(w,C)}{v(y,B)}\Bigr).
\]
This holds for every admissible witness $y$, so after taking the minimum in the definition of
$v(x,C/B)$ we conclude that $v(x,A) \leq v(x,C/B)$.

Conversely, assume $A \Sem C/B$. Let $w \in \Frag(\mathcal C)$. We show that
$v(w,A\concat B) \leq v(w,C)$. Consider any decomposition $w=xy$ with $x,y \in \Frag(\mathcal C)$.
If $v(y,B)=0$, then
\[
v(x,A)\,v(y,B)=0 \leq v(w,C).
\]
If $v(y,B)>0$, then by the assumption $A \Sem C/B$ we have
\[
v(x,A) \leq v(x,C/B).
\]
Now the definition of $v(x,C/B)$ is a minimum over admissible witnesses, one of which is
precisely the present decomposition $xy=w$. Hence
\[
v(x,C/B) \leq \min\Bigl(1,\frac{v(w,C)}{v(y,B)}\Bigr).
\]
Combining the two inequalities and multiplying by $v(y,B)$ gives
\[
v(x,A)\,v(y,B) \leq v(w,C).
\]
Thus every candidate product score for a split of $w$ is bounded above by $v(w,C)$, and so
its maximum is as well:
\[
v(w,A\concat B) \leq v(w,C).
\]
Therefore $A\concat B \Sem C$.
\end{proof}

\begin{proposition}[Left residuation]
For all formulae $A,B,C$,
\[
A\concat B \Sem C
\quad\Longleftrightarrow\quad
B \Sem A\backslash C.
\]
\end{proposition}

\begin{proof}
The proof is symmetric to the previous one. Instead of decompositions of the form $xy=w$, one
uses decompositions of the form $yx=w$, which are exactly the contexts quantified over in the
definition of $v(x,A\backslash C)$.
\end{proof}

\section{Transitivity}
\begin{theorem}[Transitivity]
If $A \Sem B$ and $B \Sem C$, then $A \Sem C$.
\end{theorem}

\begin{proof}
Let $x \in \Frag(\mathcal C)$. By assumption,
\[
v(x,A) \leq v(x,B)
\qquad\text{and}\qquad
v(x,B) \leq v(x,C).
\]
Hence, by transitivity of $\leq$ on $[0,1]$, we obtain $v(x,A) \leq v(x,C)$. Since $x$ was
arbitrary, $A \Sem C$.
\end{proof}

\section{Main soundness theorem}
\begin{theorem}[Soundness]
If $A \vdash_{\LC} B$, then $A \Sem B$.
\end{theorem}

\begin{proof}
Identity, associativity, residuation, and transitivity have all been shown valid under $\Sem$.
Since $\vdash_{\LC}$ is generated by these principles, every derivable inequality is semantically
valid. Therefore $A \vdash_{\LC} B$ implies $A \Sem B$.
\end{proof}

\begin{remark}
The proof crucially uses the cleaned-up treatment of the residuals. In the earlier draft, division
by zero and the status of unattested strings made the soundness argument unstable. Once the
domain is restricted to observed fragments and the residual clauses quantify only over positive
denominators, the argument becomes straightforward.
\end{remark}

\chapter{Future Work}

\section{Completeness for the probabilistic semantics}
The most immediate next step is a completeness theorem, or a proof that completeness fails for
this semantics. Either result would be valuable. A completeness proof would show that the
semantics captures exactly the inferential content of the associative Lambek calculus. A failure
of completeness would be equally informative, because it would identify the precise point at
which corpus-relative probabilistic information departs from the classical residuated behaviour.

The main obstacles are already visible. The semantics is finite, numerical, and tied to attested
fragments. Standard canonical constructions for classical Lambek models do not immediately
produce such corpus-relative valuations. Moreover, the use of maxima and minima in the
semantic clauses suggests that any completeness proof may need to build a model whose
numerical values are carefully engineered rather than merely order-theoretic.

\section{Extending the formal model}
Several extensions would make the semantics more realistic.
\begin{enumerate}[label=(\roman*)]
    \item \textbf{Out-of-vocabulary handling.} A smoothing method should assign principled values to
    unseen words or fragments.
    \item \textbf{Learning lexical probabilities.} Instead of assuming $\Lex$, one should estimate it from
    annotated or weakly supervised corpora.
    \item \textbf{Beyond observed fragments.} The present semantics could be extended from
    $\Frag(\mathcal C)$ to all of $\Sigma^+$, perhaps with backoff or compositional generalisation.
    \item \textbf{Normalisation over derivations.} A fuller probabilistic semantics could attach
    probabilities directly to derivation trees, not only to string--type pairs.
\end{enumerate}

\section{Implementation}
A concrete implementation would be a natural and worthwhile continuation of the thesis.
One plausible route is a chart parser or dynamic-programming parser for the associative Lambek
calculus with weighted lexical entries. The product clause already suggests a bottom-up dynamic
program: for each substring, compute the best score for each complex type by maximising over
binary splits. Residual values could then be computed from attested contexts or approximated by
statistics gathered during training.

A practical implementation project could proceed in four stages.
\begin{enumerate}[label=(\arabic*)]
    \item Build a lexicon of atomic type probabilities from an annotated corpus.
    \item Implement a weighted Lambek chart parser over substrings.
    \item Compare the parser's ambiguity-resolution behaviour with PCFG and probabilistic CCG
    baselines.
    \item Investigate whether the semantic values provide useful ranking information for downstream
    semantic interpretation.
\end{enumerate}

Such an implementation would also help test whether the present semantics is best interpreted as
(a) a genuine probabilistic semantics, (b) a scoring semantics for parse ranking, or (c) the first
step toward a fuller statistical Lambek model.

\section{Final remarks}
This thesis has aimed to do one thing carefully rather than many things incompletely: to turn a
promising but technically unstable sketch into a coherent formal proposal. The result is a
corpus-based probabilistic semantics for the associative Lambek calculus together with a
soundness theorem. That already provides a viable foundation for a larger research programme.
The next stage is clear: completeness, implementation, and empirical evaluation.

\printbibliography

\end{document}
